% part: intuitionistic-logic
% chapter: introduction
% section: constructive-reasoning
%READERNOTE{SOL6-A345: SOL6-A195 मधील गणिताने वैध पण ऐच्छिक अपरिमेयता-सिद्धता भाषांतरातून मागे घेतली. मूळाची संक्षिप्त मांडणी, निवडलेली मूल्ये आणि घाताची गणना जपली; विस्तृत सिद्धता अंमलात न आणलेली सुधारणा-सूचना म्हणून नोंदवली आहे.}

\documentclass[../../../include/open-logic-chapter]{subfiles}

\begin{document}

\olfileid{int}{int}{cr}

\section{रचनाशील तर्कविचार}

मोडल संकारक किंवा द्वितीय-क्रम संख्यापक जोडून अभिजात
तर्कशास्त्राचा विस्तार करता येतो; त्याउलट
अंतःप्रज्ञावादी तर्कशास्त्र अभिजात तर्कशास्त्रावर
मर्यादा घालते, म्हणून ते ``अभिजातेतर'' आहे.
अभिजात तर्कशास्त्र अनेक अर्थांनी \emph{अरचनाशील}
आहे. एका प्रकारच्या रचनाशील गणिताला वैशिष्ट्यपूर्ण
असणारा अधिक ``रचनाशील'' तर्कविचार अंतःप्रज्ञावादी
तर्कशास्त्रात पकडायचा आहे. यामागील काही प्रेरणा
पुढील उदाहरणांतून स्पष्ट होतील.

समजा, एखादी व्यक्ती असा दावा करते की तिने एक नैसर्गिक
संख्या~$n$ शोधली आहे: $n$ सम असेल, तर रीमान
गृहीतक सत्य आहे; आणि $n$ विषम असेल, तर ते असत्य
आहे. छान बातमी!{} रीमान गृहीतक सत्य आहे की नाही
हा गणितातील मोठ्या अनिर्णीत प्रश्नांपैकी एक आहे.
या व्यक्तीने तो प्रश्न गणनेचा, म्हणजे विशिष्ट संख्या
सम आहे की नाही हे ठरवण्याचा, प्रश्न केला असे दिसते.

मग~$n$ चे ते जादुई मूल्य कोणते? ती व्यक्ती असे
सांगते: रीमान गृहीतक सत्य असेल, तर~$n$ ही नैसर्गिक
संख्या $2$ आहे; अन्यथा ती $3$ आहे.

तुम्ही रागाने तुमचे पैसे परत मागता. अभिजात दृष्टिकोनातून
वरील वर्णन खरोखरच~$n$ चे एकमेव मूल्य ठरवते;
पण तुम्हाला हवे आहे ते~$n$ चे \emph{स्पष्टपणे}
दिलेले मूल्य.

आता आणखी एक, कदाचित कमी कृत्रिम, उदाहरण पाहू.
अपरिमेय संख्येचा परिमेय घात घेऊन परिमेय उत्तर मिळू
शकते, हे माहीत आहे. उदाहरणार्थ, $\sqrt{2}^2 = 2$.
पण अपरिमेय संख्येचा \emph{अपरिमेय} घात घेतल्यावरही
परिमेय उत्तर मिळू शकते का, हे तितके स्पष्ट नाही.
पुढील प्रमेय या प्रश्नाचे होकारार्थी उत्तर देते:

\begin{thm}
$a$ आणि $b$ अशा अपरिमेय संख्या आहेत की $a^b$ परिमेय आहे.
\end{thm}

\begin{proof}
$\sqrt{2}^{\sqrt{2}}$ पाहा. ही संख्या परिमेय असेल,
तर काम झाले: $a = b = \sqrt{2}$ घेता येईल. अन्यथा
ती अपरिमेय आहे. मग
\[
(\sqrt{2}^{\sqrt{2}})^{\sqrt{2}} = \sqrt{2}^{\sqrt{2} \cdot
  \sqrt{2}} = \sqrt{2}^2 = 2,
\]
आणि हे परिमेय आहे. म्हणून या वेळी $a$ म्हणून
$\sqrt{2}^{\sqrt{2}}$ आणि $b$ म्हणून~$\sqrt 2$ घ्या.
\end{proof}

ही सिद्धता वैध आहे का? बहुतेक गणितज्ञांना ती वैध
वाटते. तरीही इथे थोडे असमाधानकारक काहीतरी आहे:
विशिष्ट गुणधर्म असलेल्या वास्तव संख्यांच्या जोडीचे
अस्तित्व सिद्ध केले, पण ती जोडी \emph{कोणती} हे
सांगता आले नाही. हाच निष्कर्ष अशा रीतीनेही सिद्ध
करता येतो की $a$, $b$ ही जोडी सिद्धतेत \emph{दिली}
जाते: $a = \sqrt{3}$ आणि $b = \log_3 4$ घ्या. मग
\[
a^b = \sqrt{3}^{\log_3 4} = 3^{1/2 \cdot \log_3 4} = (3^{\log_3
  4})^{1/2} = 4^{1/2}= 2,
\]
कारण $3^{\log_3 x} = x$.

पहिल्या सिद्धतेतील पावलांसारखी पावले ग्राह्य नसतील
असा तर्कविचार अंतःप्रज्ञावादी तर्कशास्त्रात पकडायचा
आहे. $!A(x)$ ची पूर्तता करणारा $x$ अस्तित्वात आहे
हे सिद्ध करण्यासाठी दुसऱ्या सिद्धतेप्रमाणे एखादा
विशिष्ट~$x$ आणि तो $!A$ ची पूर्तता करतो याची
सिद्धता द्यावी लागते. $!A$ किंवा $!B$ सत्य आहे
हे सिद्ध करायचे असेल, तर या दोहोंपैकी एकाची
सिद्धता देता आली पाहिजे.

औपचारिक भाषेत सांगायचे तर, अभिजात तर्कशास्त्राची
!!a{derivation} व्यवस्था विशिष्ट प्रकारे मर्यादित
केल्यावर अंतःप्रज्ञावादी तर्कशास्त्र मिळते. गणिती
दृष्टिकोनातून या केवळ औपचारिक निगमन व्यवस्था
आहेत; पण आधी सांगितल्याप्रमाणे त्यांत एका प्रकारचा
गणिती तर्कविचार पकडायचा आहे. तो तर्कविचार
गणिताविषयीच्या विशिष्ट तात्त्विक दृष्टिकोनातून
(उदाहरणार्थ, ब्रौवर यांच्या अंतःप्रज्ञावादातून)
समर्थनीय मानता येईल; किंवा बिशप यांच्या
रचनाशीलतावादाच्या धर्तीवर अधिक ``मूर्त'' आणि
समाधानकारक गणिती तर्कविचार मानता येईल. औपचारिक
वर्णन अनौपचारिक प्रेरणा खरोखर पकडते का, यावरही
वाद करता येईल. पण आपली तात्त्विक भूमिका काहीही
असो, अंतःप्रज्ञावादी तर्कशास्त्राचा औपचारिकरीत्या
मांडलेल्या तर्कशास्त्राप्रमाणे अभ्यास करता येतो;
आणि अनेक गणिती तर्कशास्त्रज्ञांना तो अभ्यास
रसपूर्ण वाटतो.

\end{document}
